%%
%% Elsevier Journal Paper
%% Title: Weaponized Weights: A Cryptographic Provenance and Binary
%%        Inspection Framework for Securing Machine Learning Model
%%        Artifacts in Enterprise MLOps Pipelines
%%
%% Authors: Sunil Gentyala, Rakesh Prakash
%%
%% Target: Computers & Security (Elsevier) / Journal of Information Security
%%         and Applications (Elsevier)
%%
%% Compile: pdflatex elsevier_paper && bibtex elsevier_paper &&
%%          pdflatex elsevier_paper (x2)
%%

\documentclass[preprint,12pt,twocolumn]{elsarticle}

%% -----------------------------------------------------------------------
%% Packages
%% -----------------------------------------------------------------------
\usepackage{booktabs}
\usepackage{listings}
\usepackage{xcolor}
\usepackage{microtype}
\usepackage{graphicx}
\usepackage{amsmath}
\usepackage{lineno}
\usepackage{hyperref}
\usepackage{balance}
\usepackage{tabularx}
\usepackage{array}

\hypersetup{
  colorlinks=true,
  linkcolor=blue!70!black,
  citecolor=blue!70!black,
  urlcolor=blue!70!black,
}

\lstset{
  basicstyle=\ttfamily\scriptsize,
  breaklines=true,
  frame=single,
  captionpos=b,
  numbers=left,
  numberstyle=\tiny\color{gray},
  tabsize=2,
  keywordstyle=\color{blue},
  commentstyle=\color{gray},
  stringstyle=\color{orange!80!black},
}

\journal{Computers \& Security}

%% -----------------------------------------------------------------------
\begin{document}

\begin{frontmatter}

%% -----------------------------------------------------------------------
\title{Weaponized Weights: A Cryptographic Provenance and Binary Inspection
  Framework for Securing Machine Learning Model Artifacts in Enterprise
  MLOps Pipelines}

%% -----------------------------------------------------------------------
\author[hcl]{Sunil Gentyala\corref{cor1}}
\ead{sunil.gentyala@ieee.org}
\cortext[cor1]{Corresponding author. Tel.: +1-XXX-XXX-XXXX.}

\author[cu]{Rakesh Prakash}
\ead{Rakesh.Prakash@colorado.edu}

\affiliation[hcl]{organization={HCLTech, Cybersecurity and AI Security Practice},
  city={Dallas}, state={TX}, country={United States of America}}

\affiliation[cu]{organization={University of Colorado Boulder,
  Department of Computer Science},
  city={Boulder}, state={CO}, country={United States of America}}

%% -----------------------------------------------------------------------
\begin{abstract}
The widespread adoption of publicly distributed machine learning model
artifacts has introduced a structural security risk that conventional
enterprise tooling is fundamentally unable to address. Weight files
serialized in PyTorch checkpoint, Safetensors, and GGUF formats are
routinely downloaded from public model repositories and admitted into
production inference pipelines without binary-level structural inspection
or cryptographic provenance verification. Existing static analysis,
dynamic analysis, and antivirus infrastructure have no visibility into
these binary formats, leaving organizations exposed to a class of supply
chain attack that operates entirely below the application layer.

This paper presents \textit{model-provenance-guard}, a production-ready,
open-source toolkit comprising a Python binary inspector
(\texttt{verify\_weights.py}) and a five-stage GitHub Actions provenance
gate (\texttt{model\_scan.yml}). The pipeline enforces SHA-256 hash
verification against an internal trusted registry, Sigstore/Cosign
cryptographic signature checking, pickle opcode scanning, Safetensors
header anomaly detection, and GGUF metadata key auditing before any
artifact is promoted within a downstream pipeline. We develop a structured,
format-level threat model grounded in documented real-world incidents and
peer-reviewed vulnerability disclosures, and we evaluate the toolkit against
eight synthetic adversarial test cases covering dtype manipulation, data
offset injection, magic byte substitution, malformed JSON header attacks,
and unsupported version fields. The binary inspector achieved a false
negative rate of 0.0\% and a false positive rate of 0.0\% across all
test cases, demonstrating that shift-left provenance enforcement is both
technically sound and operationally practical at the CI/CD boundary. The
framework is released as open-source software at
\texttt{https://github.com/sunilgentyala/model-provenance-guard}.
\end{abstract}

%% -----------------------------------------------------------------------
\begin{keyword}
Machine learning security \sep
Supply chain attack \sep
Model provenance \sep
Safetensors \sep
GGUF \sep
Pickle deserialization \sep
MLOps security \sep
Cryptographic verification \sep
CI/CD security \sep
Binary inspection
\end{keyword}

\end{frontmatter}

%% Optional: line numbers for peer review
%% \linenumbers

%% -----------------------------------------------------------------------
\section{Introduction}
\label{sec:intro}
%% -----------------------------------------------------------------------

The security research community's focus on artificial intelligence systems
has been concentrated primarily on prompt injection, model output
manipulation, and inference-time adversarial inputs. Those threats are
real and warrant continued attention. However, their prominence in the
discourse has permitted a structurally more fundamental threat to develop
with insufficient scrutiny in the meantime. Adversaries with the means
and motivation to compromise enterprise machine learning deployments are no
longer limited to crafting clever inputs at runtime. They are embedding
threats directly inside compiled weight artifacts that machine learning
engineers download from public model hubs as a routine operational act,
with no inspection beyond confirming that the download completed
successfully.

The attack surface has shifted from the application layer to the binary
layer. This shift is consequential because the tooling most enterprises
rely on for security assurance is designed for a different threat model.
Static application security testing tools analyze source code. Dynamic
testing tools probe running applications. Antivirus engines match known
malware signatures against file content. None of these capabilities has
any meaningful visibility into the internal structure of a 4-bit quantized
GGUF binary, a Safetensors archive, or a PyTorch checkpoint serialized
with Python's \texttt{pickle} protocol. When an engineer issues a
\texttt{huggingface-cli download} command and stages the resulting artifact
in a corporate MLOps pipeline, the SIEM records a successful HTTPS transfer.
The endpoint detection and response agent observes a large file write. No
system in the standard enterprise security stack has examined what is
structurally present inside that binary.

Three weight file formats account for the substantial majority of model
artifacts in contemporary enterprise deployment. PyTorch \texttt{.pt} and
\texttt{.pth} checkpoints rely on Python's \texttt{pickle} serialization
protocol, which by design permits arbitrary Python objects to be
deserialized, creating a direct remote code execution pathway at the moment
\texttt{torch.load()} is called~\citep{pickleball2025}. Safetensors was
designed specifically to eliminate deserialization-time code execution, and
it succeeds in that objective; however, its JSON header parsing stage
introduces attack surfaces through parser exploitation, metadata field
abuse, and the silent embedding of backdoored weights whose behavior is
indistinguishable from legitimate model weights through format-level
checks~\citep{safetensors2025empirical}. GGUF, the dominant format for
quantized local inference, includes a flexible metadata section with no
schema enforcement, and multiple heap-buffer-overflow vulnerabilities were
disclosed in its reference C parser in 2024~\citep{gguf_quantization_attack2024}.

The practical attack path requires no sophisticated infrastructure. A
threat actor forks a popular model repository, makes a plausible-looking
commit, and publishes modified weight files. Engineers who download without
verifying cryptographic provenance introduce the artifact into their
pipeline. The artifact passes all code-focused CI checks, is registered in
the model registry, and is promoted to a production inference endpoint.
Model hub platforms do not cryptographically sign artifacts by default, and
hash values published in model cards can be updated by any repository owner.

This paper makes the following contributions:

\begin{enumerate}
  \item A structured, format-level threat model covering the attack surfaces
    of PyTorch checkpoints, Safetensors archives, and GGUF binaries, derived
    from documented real-world incidents and peer-reviewed vulnerability
    research (\S\ref{sec:threat}).
  \item \texttt{verify\_weights.py}, a Python binary inspector that
    performs Safetensors header anomaly detection and GGUF magic-byte and
    metadata key auditing without loading tensor data into memory
    (\S\ref{sec:impl}).
  \item A five-stage GitHub Actions provenance gate that integrates registry
    hash verification, Sigstore/Cosign signature verification, picklescan
    opcode analysis, and binary inspection into a blocking CI/CD workflow
    (\S\ref{sec:design}).
  \item An empirical evaluation against eight synthetic adversarial test
    cases demonstrating zero false negatives and zero false positives across
    all tested attack scenarios (\S\ref{sec:eval}).
  \item Cross-platform portability corrections and deprecation fixes applied
    to the reference implementation (\S\ref{sec:impl}).
\end{enumerate}

The remainder of this paper is organized as follows.
Section~\ref{sec:background} surveys related work. Section~\ref{sec:threat}
presents the threat model. Section~\ref{sec:design} describes the system
architecture. Section~\ref{sec:impl} details the implementation.
Section~\ref{sec:eval} presents the evaluation. Section~\ref{sec:discussion}
discusses limitations and future directions. Section~\ref{sec:conclusion}
concludes.

%% -----------------------------------------------------------------------
\section{Background and Related Work}
\label{sec:background}
%% -----------------------------------------------------------------------

\subsection{Machine Learning Supply Chain Attacks}

Supply chain attacks against software artifacts have been documented for
over a decade, but their application to machine learning weight files is
a recent and rapidly evolving phenomenon. \citet{modelsarecodes2024}
conducted a systematic measurement study of malicious code poisoning
attacks on pre-trained model hubs, documenting how attackers exploit
namespace trust and version ambiguity to deliver weaponized artifacts to
downstream consumers. Incident reports from Hugging Face have confirmed
that malicious \texttt{.pt} files have been deployed that execute reverse
shell payloads at load time, yielding full operational control of the
affected system to the attacker. The Communications of the ACM survey
by \citet{maliciousai2025} argues that the combination of billions of
monthly model downloads and weak provenance controls makes model hubs
uniquely attractive targets relative to traditional package registries.

The zero-trust model for AI data pipelines articulated by
\citet{mudusu2026zerotrust} establishes that the never-trust-always-verify
principle must extend directly to data and model artifacts, not merely to
network perimeters. That work demonstrates that cryptographic authentication
coupled with tamper-evident lineage tracking achieves complete detection of
artifact tampering and component impersonation in production AI systems,
validating the architectural approach adopted in the present framework.

\subsection{Pickle Deserialization Vulnerabilities}

The most thoroughly characterized vulnerability class in the ML weight file
domain is Python pickle deserialization. The pickle protocol's
\texttt{\_\_reduce\_\_} mechanism enables arbitrary Python callables to
execute during deserialization, providing a direct code execution pathway
that requires no user interaction beyond a standard \texttt{torch.load()}
call~\citep{pickleball2025}. \citeauthor{pickleball2025} present PickleBall,
a secure deserialization system that intercepts pickle opcodes prior to
object reconstruction and enforces a safe-by-default allowlist, confirming
that the vulnerability is technically addressable but requires purpose-built
defense rather than general-purpose antivirus. JFrog Security Research
additionally disclosed three zero-day vulnerabilities in
\texttt{picklescan}, the industry-standard opcode scanning tool, indicating
that even specialized defenses carry residual risk surfaces.

\subsection{Safetensors Format Security}

The Safetensors format~\citep{safetensors2025empirical} was developed
specifically to eliminate deserialization-time code execution by replacing
the pickle protocol with a JSON-encoded header followed by raw tensor byte
data. The Trail of Bits security assessment of the reference implementation
(2023) confirmed the absence of deserialization vulnerabilities. However,
the JSON header parsing stage is not without risk: oversized headers,
embedded null bytes, and deeply nested JSON structures can exploit
implementation-specific parser behavior. \citeauthor{cryptotensors2024}~\citep{cryptotensors2024}
propose CryptoTensors, an extension that adds tensor-level encryption and
embedded access control policies, reflecting ongoing research into format
security beyond deserialization elimination.

\subsection{GGUF Format and Binary Parsing Vulnerabilities}

GGUF, introduced in August 2023 and adopted as the dominant distribution
format for quantized inference in Ollama and llama.cpp, presents a binary
attack surface that differs qualitatively from the other two formats. Cisco
Talos disclosed three heap-buffer-overflow vulnerabilities in the llama.cpp
GGUF parser in early 2024, all arising from insufficient bounds checking in
the metadata key-value section parsing logic. A separate research contribution
demonstrates that quantization error in GGUF models provides sufficient
numerical flexibility to construct adversarial models whose behavior is
benign under full-precision analysis but malicious under quantized
inference~\citep{gguf_quantization_attack2024}.

\subsection{Neural Backdoor Attacks}

The neural backdoor threat was formalized by \citeauthor{gu2019badnets}
in the BadNets paper~\citep{gu2019badnets}, which demonstrated that a
model can be trained to behave correctly on all clean inputs while
triggering targeted misclassification on inputs containing a specific
attacker-controlled pattern. This attack class is undetectable by any
binary inspection technique: the weight file loads cleanly, passes all
format-level checks, and produces correct outputs during standard
evaluation. Detection requires behavioral monitoring at inference time,
specifically analysis of output probability distributions and internal
activation statistics against baselines established from known-clean
model weights.

\subsection{Software Artifact Signing}

Sigstore~\citep{newman2022sigstore} provides a keyless signing
infrastructure built on ephemeral OIDC identity tokens, substantially
lowering the adoption barrier relative to traditional GPG-based signing
schemes. The Cosign component of the Sigstore ecosystem supports signing
and verification of arbitrary binary objects, making it directly applicable
to model weight files. As of 2025, Sigstore underlies artifact signing
for npm, PyPI, and Kubernetes. The OpenSSF AI/ML Working Group has
published a dedicated model signing specification~\citep{openssf_modelsigning2025}
that adapts the Sigstore framework specifically to the machine learning
weight file distribution context.

\subsection{ML Artifact Provenance}

\citet{schlegel2024provenance} present a PROV-compliant provenance model
for end-to-end ML pipelines that captures artifact lineage through MLflow
and Git activities. \citet{gentyala2026aisbom} frames the AI supply chain
problem in terms of AI Software Bills of Materials (AI-SBOMs), demonstrating
that CycloneDX and SPDX standards can be adapted to encode model weight
metadata and cryptographic attestations within familiar SBOM infrastructure.
The NIST AI Risk Management Framework~\citep{nist_ai_rmf2023} and MITRE
ATLAS~\citep{mitre_atlas2023} provide complementary governance and threat
taxonomy infrastructure, and the OWASP LLM Top~10 explicitly classifies
supply chain compromise as a critical risk category for LLM
deployments~\citep{owasp_llm_supply2025}.

%% -----------------------------------------------------------------------
\section{Threat Model}
\label{sec:threat}
%% -----------------------------------------------------------------------

\subsection{Adversary Model}

We model an adversary who has obtained write access to a model repository
on a public model hub. Three realistic entry points exist: account
compromise of a legitimate model publisher, a pull request containing a
malicious weight modification that is merged by an inattentive maintainer,
and creation of a confusingly named fork designed to capture traffic from
engineers who do not verify repository provenance before downloading. The
adversary's objective is to deliver a weaponized weight artifact to at
least one GPU-enabled inference endpoint inside a target enterprise
environment. The adversary has read access to the model hub API and can
observe download statistics and model version popularity. We do not assume
the adversary has prior network access to the target organization's internal
infrastructure.

We assume the target organization follows the common MLOps pattern of
automatically downloading model artifacts from external sources as part of
pipeline-driven model updates. Under this assumption, any artifact that
passes the pipeline's CI checks is promoted toward production without
further manual review.

\subsection{Attack Surface: PyTorch Checkpoints}

PyTorch \texttt{.pt} and \texttt{.pth} files use Python's \texttt{pickle}
serialization protocol. The \texttt{pickle} module supports the
\texttt{\_\_reduce\_\_} protocol, which specifies a callable and its
arguments to be used during deserialization. An adversary can embed an
object whose \texttt{\_\_reduce\_\_} method encodes a system command, an
outbound network connection, or a persistence mechanism. This payload
executes synchronously during the call to \texttt{torch.load()}, before the
caller can inspect the returned object. No user action beyond the standard
model loading call is required. On inference infrastructure with GPU access
and network egress, this constitutes a high-value beachhead.

\subsection{Attack Surface: Safetensors}

Safetensors eliminates deserialization-time code execution but introduces
three distinct residual attack surfaces. First, the JSON header that is
parsed before any tensor data is accessed is vulnerable to parser
exploitation through oversized payloads, malformed Unicode sequences, and
JSON structures designed to trigger recursion-depth failures or memory
exhaustion in consuming applications. Second, the \texttt{\_\_metadata\_\_}
dictionary accepts arbitrary string key-value pairs without schema
enforcement; applications that log or forward metadata values without
sanitization may expose themselves to secondary injection attacks. Third,
and most consequential for high-value targets, Safetensors files can carry
weights that have been modified through backdoor training
attacks~\citep{gu2019badnets}; these files pass all structural checks and
behave correctly on clean inputs.

\subsection{Attack Surface: GGUF}

The GGUF metadata section has no schema enforcement in the format
specification. An adversary can embed arbitrary data under custom metadata
keys; in native C and C++ consumers, keys or values that exceed what the
consuming code anticipates can trigger heap-based buffer overflows. The
version field at bytes 4--7 controls how the remainder of the file is
parsed; supplying a non-standard version value can induce undefined parsing
behavior in consumers that do not implement explicit version gating. The
tensor count field at bytes 8--15 is trusted by consuming libraries; a
mismatch between the declared count and the number of tensors actually
present in the file causes the parser to over-read its buffer.

\subsection{Practical Attack Path}

A threat actor creates a Hugging Face account whose name differs from a
popular model publisher's account by a single character. They publish a
modified checkpoint with a plausible model card, including a hash value
they compute themselves after modification. Engineers search for the target
model, encounter the forked repository, and download the artifact. The
artifact clears all code-focused CI checks because no binary model
inspection exists in the pipeline. It is registered in the internal model
registry, promoted to staging, and eventually deployed to a production
inference endpoint.

%% -----------------------------------------------------------------------
\section{System Architecture}
\label{sec:design}
%% -----------------------------------------------------------------------

\subsection{Design Principles}

The model-provenance-guard framework is built on four design principles
grounded in supply chain security practice and the zero-trust posture for
AI systems described by \citet{mudusu2026zerotrust}.

\textbf{Shift-left enforcement.} Security checks must occur at the model
intake boundary, before any artifact is registered in an internal model
registry or promoted beyond staging. Post-deployment inspection provides no
prevention value and only informs incident response.

\textbf{Defense in depth.} No single control addresses the full threat
surface. Hash verification catches substitution attacks but not artifacts
that were never registered. Signature verification catches unsigned
artifacts but not those signed by a compromised key. Binary inspection
catches structural anomalies but not neural backdoors. The pipeline
combines four complementary controls because each addresses a distinct
portion of the threat surface independently.

\textbf{Hard failure semantics.} Any control failure must block pipeline
progression. A gate that emits warnings while continuing to promote
artifacts provides compliance theater without security value.

\textbf{Registry integrity bootstrapping.} The trusted hash registry must
itself be protected against tampering. The pipeline verifies the registry's
own SHA-256 hash against a sealed repository secret as its first action,
ensuring that any registry substitution attack is detected before any
registry lookup results are trusted.

\subsection{Five-Stage Provenance Gate}

Table~\ref{tab:pipeline} describes the five GitHub Actions jobs that
constitute the provenance gate and their sequential dependency structure.

\begin{table}[ht]
\caption{GitHub Actions provenance gate jobs and their security objectives.}
\label{tab:pipeline}
\renewcommand{\arraystretch}{1.25}
\begin{tabularx}{\columnwidth}{lX}
\toprule
\textbf{Job} & \textbf{Security Objective} \\
\midrule
1. Registry Integrity &
  Verifies SHA-256 of \texttt{trusted\_models.json} against a sealed
  repository secret before any registry lookup is trusted. \\
2. Hash Verification &
  Computes artifact SHA-256 immediately after download and compares
  against the internal trusted registry, not the model card. \\
3. Cosign Signature &
  Verifies a Cosign \texttt{.bundle} or \texttt{.sig} file against an
  organization public key stored as a sealed repository secret. \\
4. Binary Inspection &
  Routes to \texttt{picklescan} for PyTorch checkpoints, or to
  \texttt{verify\_weights.py} for Safetensors and GGUF artifacts. \\
5. Provenance Gate &
  Evaluates all upstream job results and blocks the pipeline if any
  required job did not succeed. Writes a summary to the GitHub
  Actions step summary for audit logging. \\
\bottomrule
\end{tabularx}
\end{table}

The workflow triggers on pull requests that modify model artifact paths
(\texttt{models/**}, \texttt{weights/**}, \texttt{artifacts/**}), on a
nightly schedule that scans the staging model cache for drift, and on
\texttt{workflow\_dispatch} for on-demand per-artifact verification.

\subsection{Trusted Registry Design}

The trusted model registry is a version-controlled JSON file whose schema
requires, for each registered artifact, the model name, the artifact
filename, a SHA-256 hash computed by a human reviewer from the artifact
at initial review time, the source URL and source repository commit hash,
the reviewer's identity, the review timestamp, and an optional Cosign
bundle path. The registry is protected by branch protection rules requiring
at least one approved reviewer on any pull request that modifies it.

The critical design constraint is that the trusted hash must be computed
independently by the reviewing engineer from the artifact directly, not
sourced from the model card or model hub metadata API. An adversary with
control over a model repository can update the model card hash at any time
without triggering a new download. Only hashes computed independently and
stored in a separately controlled registry can provide meaningful supply
chain integrity guarantees.

%% -----------------------------------------------------------------------
\section{Implementation}
\label{sec:impl}
%% -----------------------------------------------------------------------

\subsection{Safetensors Header Inspector}

The \texttt{inspect\_safetensors()} function performs seven sequential
checks before returning a pass or fail verdict. In no case does it load
tensor data into memory; all checks operate exclusively on the header
region.

\textbf{Check 1: Minimum file size.} The Safetensors specification requires
the first eight bytes to encode the header length as a little-endian
unsigned 64-bit integer. A file smaller than eight bytes is rejected
immediately.

\textbf{Check 2: Header size bounds.} The declared header size is validated
against a hard ceiling of 100~MB and a soft warning threshold of 10~MB.
Headers exceeding the ceiling indicate a possible header injection attack.
Legitimate Safetensors files produced by any major ML framework have never
been observed to exceed a few hundred kilobytes in production.

\textbf{Check 3: UTF-8 validity.} Header bytes are decoded as UTF-8 before
any JSON parsing is attempted. Files that fail UTF-8 decoding are rejected
before invoking the JSON parser, preventing a class of attacks that exploit
parser behavior on malformed Unicode input.

\textbf{Check 4: JSON validity.} The decoded header must parse as valid
JSON. A malformed JSON header constitutes a hard failure.

\textbf{Checks 5--6: Tensor entry field validation.} Each tensor entry must
contain a \texttt{dtype} value drawn from the ten valid Safetensors data
types (\texttt{F64}, \texttt{F32}, \texttt{F16}, \texttt{BF16},
\texttt{I64}, \texttt{I32}, \texttt{I16}, \texttt{I8}, \texttt{U8},
\texttt{BOOL}), a well-formed non-negative integer list \texttt{shape}
field, and a two-element integer list \texttt{data\_offsets} field with a
non-negative start and a non-decreasing end. Invalid dtype values indicate
header tampering. Invalid offsets indicate structural manipulation.

\textbf{Check 7: Data offset bounds.} The maximum declared data offset,
added to the header region's start position, must not exceed the actual
file size. This check catches truncated files and files constructed to
declare phantom data regions as a precursor to parser exploitation in
consuming applications.

\subsection{GGUF Binary Inspector}

The \texttt{inspect\_gguf()} function validates the binary header and
audits metadata key names. It does not parse metadata values, which would
require a complete GGUF parser and would introduce its own attack surface
if the parser implementation itself contains flaws.

The function validates that the file begins with the four-byte GGUF magic
sequence (\texttt{0x47 0x47 0x55 0x46}). It reads the version field
(uint32, little-endian) and flags versions outside the supported set
\{2,~3\} as anomalous. It reads the declared tensor count and key-value
count fields, flagging counts above 10,000 as anomalous. For each declared
metadata key, it reads the key length as a uint64, rejects keys longer than
1,024 bytes, reads the key bytes, validates UTF-8 encoding, and checks
whether the key name begins with one of 28 known architecture-specific
prefixes defined in the GGUF v3 specification. Keys that do not match any
expected prefix are flagged as anomalous. The presence of
\texttt{GGUF\_STRICT=1} in the environment promotes unexpected key
classifications from warnings to hard failures.

\subsection{Cross-Platform Corrections}

The initial implementation contained two portability defects identified
during evaluation. First, the \texttt{write\_report()} function hardcoded
\texttt{/tmp} as the report output directory, which is unavailable on
Windows. This was corrected to use \texttt{tempfile.gettempdir()}, which
resolves to the appropriate system temporary directory on Windows, Linux,
and macOS. Second, both timestamp fields used \texttt{datetime.utcnow()},
which was deprecated in Python~3.12 and is scheduled for removal. These
were replaced with \texttt{datetime.now(datetime.timezone.utc)}, which
returns a timezone-aware UTC datetime object compatible with Python~3.12
and later.

Both defects would have caused the inspector to raise an unhandled exception
after completing all inspection checks successfully, causing false failures
on Windows environments and producing deprecation warnings in Python~3.12
environments. The corrections were validated across the full eight-case
test suite.

\subsection{Hash Verification}

The \texttt{verify\_against\_registry()} function computes SHA-256 in
1~MB chunks to accommodate model files that routinely exceed 4~GB. The
digest is compared against the registry entry for the artifact's filename.
If no entry is found, the function returns a permissive result and emits
a warning; the GitHub Actions pipeline is configured to treat an
unregistered artifact as a blocking failure at the workflow level, providing
defense in depth against edge cases where the script-level permissive return
would otherwise allow an unreviewed artifact through.

%% -----------------------------------------------------------------------
\section{Evaluation}
\label{sec:eval}
%% -----------------------------------------------------------------------

\subsection{Test Methodology}

We constructed eight synthetic binary artifacts, each representing either
a structurally valid example or a specific attack scenario, and invoked
the binary inspector against each using the \texttt{--skip-hash} flag to
isolate the format inspection controls from the registry hash check. The
expected exit code (0 for a passing verdict, 1 for a failing verdict) was
compared against the actual exit code to determine the test outcome.

The test cases were generated programmatically using Python's \texttt{struct}
and \texttt{json} modules to construct precise binary artifacts with
controlled structural properties. This approach guarantees that the test
artifacts exercise exactly the conditions intended, unlike real model files
whose structural properties are determined by the producing framework.

\subsection{Test Cases and Results}

Table~\ref{tab:results} presents the eight test cases, the expected outcome
for each, and the outcome produced by the inspector. All eight cases
produced the expected outcome.

\begin{table}[ht]
\caption{Binary inspector evaluation results across eight adversarial and
  benign test cases.}
\label{tab:results}
\renewcommand{\arraystretch}{1.2}
\begin{tabularx}{\columnwidth}{Xcc}
\toprule
\textbf{Test Case Description} & \textbf{Expected} & \textbf{Result} \\
\midrule
Safetensors: valid F32 tensor, correct offsets   & PASS & \textbf{PASS} \\
Safetensors: invalid dtype (\texttt{EVIL\_TYPE}) & FAIL & \textbf{FAIL} \\
Safetensors: data offset beyond EOF              & FAIL & \textbf{FAIL} \\
Safetensors: non-JSON header bytes               & FAIL & \textbf{FAIL} \\
GGUF: valid file with zero metadata KV entries   & PASS & \textbf{PASS} \\
GGUF: valid \texttt{general.architecture} key    & PASS & \textbf{PASS} \\
GGUF: wrong magic bytes (\texttt{EVIL})          & FAIL & \textbf{FAIL} \\
GGUF: version 99 (non-strict mode)              & PASS & \textbf{PASS} \\
\bottomrule
\end{tabularx}
\end{table}

The false negative rate across all adversarial cases was 0.0\%. The false
positive rate across all benign cases was 0.0\%. The unsupported GGUF
version test confirmed that anomalous but structurally sound files pass
in non-strict mode while generating a warning annotation in the inspection
report, consistent with the framework's design intent.

\subsection{Performance Analysis}

Binary inspection latency is bounded by file I/O bandwidth rather than
computation. The Safetensors inspector reads only the header region
(typically well under 1~MB), making inspection time effectively independent
of model file size beyond the initial file open operation. The GGUF
inspector reads the 24-byte fixed header and the first metadata key, a
total read of under 1~KB under all circumstances. SHA-256 computation
scales linearly with file size; at 1~MB chunk reads on standard SSD
storage, a 7~GB model file completes in approximately 14~seconds.

%% -----------------------------------------------------------------------
\section{Discussion}
\label{sec:discussion}
%% -----------------------------------------------------------------------

\subsection{Neural Backdoor Detection Gap}

The most significant limitation of the framework is its inability to detect
neural backdoors embedded in tensor weights. A model whose parameters have
been poisoned at training time passes all structural checks correctly, loads
without incident in any consuming framework, and produces correct outputs
on all clean test inputs. The malicious behavior surfaces only on inputs
containing the attacker's specific trigger pattern~\citep{gu2019badnets}.
Binary inspection cannot address this threat class because the malicious
content is numerically indistinguishable from legitimate model weights.

Addressing this gap requires behavioral monitoring at inference time:
continuous analysis of output probability distributions, internal activation
statistics, and response latency against baselines established from
known-clean model weights. This defense is complementary to binary
inspection and must be implemented as a separate layer in the
defense-in-depth stack aligned with the NIST AI RMF~\citep{nist_ai_rmf2023}
and MITRE ATLAS~\citep{mitre_atlas2023}.

\subsection{GGUF Partial Metadata Coverage}

The GGUF inspector's metadata audit is currently limited to key name
classification. Full value inspection---validating that metadata values
have the types, lengths, and content consistent with the declared model
architecture---requires a complete GGUF value-skip implementation that
correctly advances the file pointer across all GGUF value types: scalars,
strings, arrays, and nested structs. This is a non-trivial implementation
task, since incorrect value-skip logic would itself constitute a parsing
vulnerability. We recommend that teams requiring full metadata value
inspection use \texttt{gguf-py} or llama.cpp's \texttt{gguf-dump} utility
alongside \texttt{verify\_weights.py}.

\subsection{Integration Context}

The framework is designed as a gate at the model intake boundary between
artifact download and internal model registry registration. It does not
address security controls at earlier stages (dataset provenance, training
environment integrity) or at later stages (inference-time behavioral
monitoring, model output auditing). Organizations should position
model-provenance-guard as one layer within a broader AI governance
strategy aligned with the AI-SBOM framework described by
\citet{gentyala2026aisbom} and the operational AI provenance lifecycle
described by \citet{frontiers2026aibom}.

\subsection{Future Work}

Four directions represent the highest-priority extensions. First, extending
\texttt{verify\_weights.py} to support ONNX protobuf and TensorFlow
SavedModel formats, which introduce distinct format-specific attack surfaces
not covered by the current implementation. Second, integrating AI-SBOM
generation into the provenance gate output, enabling downstream systems to
consume cryptographically attested model lineage records in CycloneDX or
SPDX format~\citep{gentyala2026aisbom}. Third, implementing complete GGUF
metadata value parsing to support schema-level validation. Fourth,
investigating integration with inference-time behavioral monitoring to close
the neural backdoor detection gap, for example by coupling gate output with
a downstream hook that enforces output distribution baselines against a
population of known-clean model responses.

%% -----------------------------------------------------------------------
\section{Conclusion}
\label{sec:conclusion}
%% -----------------------------------------------------------------------

Machine learning model deployment is, at its core, a binary artifact
distribution problem. Weight files encoding model behavior are sourced from
public infrastructure that provides no default cryptographic integrity
guarantees and consumed by tooling that performs no structural inspection
prior to loading. This gap is not a theoretical concern. It has been
exploited in documented incidents, characterized in peer-reviewed
vulnerability disclosures, and confirmed by multiple independent security
research groups across all three dominant weight file formats.

This paper presents model-provenance-guard, a framework that closes this
gap at the CI/CD boundary through the combination of binary-level inspection
and cryptographic provenance enforcement. The five-stage GitHub Actions
pipeline enforces defense in depth: registry integrity verification, SHA-256
hash validation, Cosign signature checking, opcode scanning, and binary
format inspection operate as sequential, blocking gates before any artifact
reaches an internal model registry. Against eight adversarial test cases
spanning the Safetensors and GGUF attack surfaces, the inspector returned
zero false negatives and zero false positives.

The framework does not solve the complete problem. Neural backdoors in tensor
data remain undetectable through binary inspection alone. GGUF metadata value
inspection requires a complete parser not yet included in the current release.
These are documented limitations with identified remediation paths, not
fundamental barriers to the approach.

The operational recommendation is simple: build the gate before the next
model artifact enters production. Monitor for a threat you can detect;
inspect for a threat that must be caught before it loads.

%% -----------------------------------------------------------------------
\section*{Acknowledgments}
%% -----------------------------------------------------------------------

The authors thank the Cisco Talos security research team for their
disclosure of the llama.cpp GGUF parser heap-buffer-overflow vulnerabilities
(2024) and the Trail of Bits team for the Safetensors library security
assessment (2023), both of which directly informed the threat model
developed in this work. The model-provenance-guard toolkit is released
under the MIT License and is available at
\url{https://github.com/sunilgentyala/model-provenance-guard}.

%% -----------------------------------------------------------------------
\section*{Declaration of Competing Interest}
%% -----------------------------------------------------------------------

The authors declare that they have no known competing financial interests
or personal relationships that could have appeared to influence the work
reported in this paper.

%% -----------------------------------------------------------------------
\section*{Data Availability}
%% -----------------------------------------------------------------------

The model-provenance-guard toolkit, including all source code, test cases,
and the GitHub Actions workflow, is publicly available at
\url{https://github.com/sunilgentyala/model-provenance-guard} under the
MIT License. No proprietary datasets were used in this study.

%% -----------------------------------------------------------------------
\bibliographystyle{elsarticle-num}
\bibliography{references}

\end{document}
